\subsection{Local Rings}

\begin{definition} \label{def:local-ring} \hlt{Local Rings}\\
A commutative ring $R$ with identity is \hlt{local} if it has a unique maximal ideal $\mathfrak{m}$. This unique ideal is the \hlt{maximal ideal} of $R$.
\end{definition}

\begin{definition} \label{def:localized-prime-ideal} \hlt{Localized Prime Ideal}\\
Let $R$ be a commutative ring with identity and let $\mathfrak{p}\in\operatorname{Spec}R$. The \hlt{localized prime ideal} is $\mathfrak{p}R_{\mathfrak{p}}\defeq(R\setminus\mathfrak{p})^{-1}\mathfrak{p}$.
\end{definition}

\begin{proposition} \label{prop:prime-localization-element-criterion}
\hlt{Elements of a Prime Localization}\\
Let $R$ be a commutative ring with identity, $\mathfrak{p}\in\operatorname{Spec}R$. Let $(a,s)\in R\times(R\setminus\mathfrak{p})$.
\begin{enumerate}[label=(\roman*)]
\setlength{\itemsep}{0pt}
\item $a/s\in\mathfrak{p}R_{\mathfrak{p}}\Longleftrightarrow a\in\mathfrak{p}$.
\item $a/s\notin\mathfrak{p}R_{\mathfrak{p}}\Longrightarrow a\in R\setminus\mathfrak{p}$.
\item If $a\in R\setminus\mathfrak{p}$, then $s/a$ is a multiplicative inverse of $a/s$ in $R_{\mathfrak{p}}$.
\end{enumerate}
\begin{proof}
\begin{enumerate}[label=(\roman*)]
\setlength{\itemsep}{0pt}
\item[$(i)$] ($\Rightarrow$) Assume $a/s\in\mathfrak{p}R_{\mathfrak{p}}$. Choose $u\in R\setminus\mathfrak{p}$ with $ua\in\mathfrak{p}$ by Proposition \ref{prop:localization-extension-contraction-calculus}(ii);  $a\in\mathfrak{p}$ by $u\notin\mathfrak{p}$ and primeness.
\item[$(i)$] ($\Leftarrow$) Assume $a\in\mathfrak{p}$. Then $a/s\in\mathfrak{p}R_{\mathfrak{p}}$ by Definition \ref{def:localized-prime-ideal}.
\item[$(ii)$] The implication holds by the contrapositive of (i).
\item[$(iii)$] The membership $s/a\in R_{\mathfrak{p}}$ holds by $a\in R\setminus\mathfrak{p}$. The identity $(a/s)(s/a)=1_{R_{\mathfrak{p}}}$ holds by Definition \ref{def:ring-localization}(ii),(iv).
\end{enumerate}
\hlt{Structural Commentary}: Extension tests numerators; swapping fractions constructs inverses.\\
\hlt{Key Insight}: Numerators outside $\mathfrak{p}$ become denominators for explicit inverses.
\end{proof}
\end{proposition}

\begin{theorem} \label{thm:prime-localization-is-local} \hlt{Prime Localizations Are Local}\\
Let $R$ be a commutative ring with identity and let $\mathfrak{p}\in\operatorname{Spec}R$.
\begin{enumerate}[label=(\roman*)]
\setlength{\itemsep}{0pt}
\item The ideal $\mathfrak{p}R_{\mathfrak{p}}$ is the unique maximal ideal of $R_{\mathfrak{p}}$.
\item The ring $R_{\mathfrak{p}}$ is local.
\item The quotient ring $R_{\mathfrak{p}}/\mathfrak{p}R_{\mathfrak{p}}$ is a field.
\end{enumerate}
\begin{proof}
\begin{enumerate}[label=(\roman*)]
\setlength{\itemsep}{0pt}
\item The relation $\mathfrak{p}R_{\mathfrak{p}}\trianglelefteq R_{\mathfrak{p}}$ holds by Proposition \ref{prop:localization-extension-contraction-calculus}(i). The ideal is proper by $\mathfrak{p}\cap(R\setminus\mathfrak{p})=\emptyset$ and Proposition \ref{prop:localization-extension-contraction-calculus}(iv).\\
Fix a proper ideal $\mathcal{J}\trianglelefteq R_{\mathfrak{p}}$ and assume $\mathcal{J}\nsubseteq\mathfrak{p}R_{\mathfrak{p}}$. Choose $x\in\mathcal{J}\setminus\mathfrak{p}R_{\mathfrak{p}}$ by the failed inclusion.\\
Choose $(a,s)\in R\times(R\setminus\mathfrak{p})$ with $x=a/s$ by Definition \ref{def:ring-localization}(iii); $a\in R\setminus\mathfrak{p}$ and $(a/s)(s/a)=1_{R_{\mathfrak{p}}}$ by Proposition \ref{prop:prime-localization-element-criterion}(ii),(iii).\\
The membership $1_{R_{\mathfrak{p}}}=x(s/a)\in\mathcal{J}$ contradicts properness by ideal absorption. Thus $\mathcal{J}\subseteq\mathfrak{p}R_{\mathfrak{p}}$.\\
Fix $\mathcal{K}\trianglelefteq R_{\mathfrak{p}}$ with $\mathfrak{p}R_{\mathfrak{p}}\subseteq\mathcal{K}$ and $\mathcal{K}\neq R_{\mathfrak{p}}$. Then $\mathcal{K}\subseteq\mathfrak{p}R_{\mathfrak{p}}$ by the proper-ideal containment.\\
Hence $\mathcal{K}=\mathfrak{p}R_{\mathfrak{p}}$, so $\mathfrak{p}R_{\mathfrak{p}}$ is maximal by Definition \ref{def:maximal-ideal}. If $\mathfrak{m}\trianglelefteq R_{\mathfrak{p}}$ is maximal, then $\mathfrak{m}\subseteq\mathfrak{p}R_{\mathfrak{p}}$ by the proper-ideal containment.\\
Therefore $\mathfrak{m}=\mathfrak{p}R_{\mathfrak{p}}$ by maximality of $\mathfrak{m}$ and properness of $\mathfrak{p}R_{\mathfrak{p}}$.
\item The ring $R_{\mathfrak{p}}$ is local by (i) and Definition \ref{def:local-ring}.
\item The quotient is a field by (i) and Theorem \ref{thm:quotient-by-maximal-ideal-is-field}.
\end{enumerate}
\hlt{Structural Commentary}: Explicit inverses force proper ideals into $\mathfrak{p}R_{\mathfrak{p}}$; quotient theory finishes.\\
\hlt{Key Insight}: Fractions outside $\mathfrak{p}R_{\mathfrak{p}}$ invert by swapping numerator and denominator.
\end{proof}
\end{theorem}

\begin{definition} \label{def:residue-fields} \hlt{Residue Fields}\\
Let $R$ be a commutative ring with identity.
\begin{enumerate}[label=(\roman*)]
\setlength{\itemsep}{0pt}
\item If $R$ is local with maximal ideal $\mathfrak{m}$, its \hlt{residue field} is $R/\mathfrak{m}$.
\item For $\mathfrak{p}\in\operatorname{Spec}R$, the \hlt{residue field at $\mathfrak{p}$} is $k(\mathfrak{p})\defeq R_{\mathfrak{p}}/\mathfrak{p}R_{\mathfrak{p}}$.
\end{enumerate}
\end{definition}

% TODO(structure): After the general commutative-ring definition of units is
% promoted from its integral-domain-only owner, add the equivalence among a unique
% maximal ideal, the ideal of nonunits, and the criterion that r or 1-r is a unit.
